\section{Endpoint Insufficiency}
\label{app:proof}

Let $\mathcal{R}_a(s)$ denote expected downstream readiness after prefix $s$ of action $a$.
Conditioned on a request arriving before the action ends, its interruption value is
\begin{equation}
I_f(a)=\int_0^{L_a} f(s)\mathcal{R}_a(s)\,ds.
\end{equation}

\paragraph{Proposition 1 (formal statement).}
For any two actions $a,b$ of equal duration $L$ such that
$\mathcal{R}_a(L)=\mathcal{R}_b(L)$ but
$\mathcal{R}_a(s)\neq\mathcal{R}_b(s)$ on a set of nonzero measure,
there exists a valid arrival density $f$ for which $I_f(a)\neq I_f(b)$.
If the sign of $\mathcal{R}_a(s)-\mathcal{R}_b(s)$ takes both positive and negative values on sets of nonzero measure, then there exist densities $f_1,f_2$ that induce opposite action rankings.

\paragraph{Proof.}
Let $D_+=\{s:\mathcal{R}_a(s)>\mathcal{R}_b(s)\}$.
If $D_+$ has nonzero measure, choose a nonnegative density whose support is contained in $D_+$.
Then the integral of the readiness difference is positive and $I_f(a)>I_f(b)$.
The symmetric construction on
$D_-=\{s:\mathcal{R}_a(s)<\mathcal{R}_b(s)\}$
gives the opposite ranking whenever $D_-$ has nonzero measure.
If only one of the sets has nonzero measure, an unequal ranking still follows.
The shared terminal value is irrelevant because a point has zero measure for a continuous arrival density; the discrete-time result follows by placing probability mass on a differing prefix.
\hfill$\square$

This proposition does not state that coverage is always a poor heuristic.
It states that no score depending only on terminal memory can preserve all interruption-time preferences when prefix readiness differs.

\section{Expanded Related Work}
\label{app:related}

\paragraph{Fixed-budget exploration.}
Classical frontier exploration, Active Neural SLAM, intrinsic-motivation methods, and recent foundation-model explorers optimize geometric coverage, learned novelty, spatial-belief quality, or endpoint checkpoint coverage
\citep{yamauchi1997frontier,chaplot2020ans,sekar2020plan2explore,ye2026look,zhang2026theory}.
AutoContext, ABot-Explorer, LMEE, and RoboEXP make pre-task experience reusable through structured context or memory
\citep{cai2026autocontext,chen2026abot,wang2026lmee,jiang2024roboexp}.
\taskname{} does not claim reusable memory as novel; it asks which evidence should be acquired first when the endpoint is not guaranteed to be reached.

\paragraph{Known-task active perception.}
Embodied QA, ActiveVOO, and ObsGraph acquire observations that support a given question or plan
\citep{das2018eqa,liu2025activevoo,lee2026obsgraph}.
They motivate evidence-sensitive action choice but can condition it on the realized request.
In \taskname{}, only a service catalog distribution is visible before interruption, so known-task systems are diagnostic upper bounds.

\paragraph{Unknown future reward.}
Reward-free RL, world-model exploration, and successor features seek knowledge reusable under later rewards
\citep{jin2020rewardfree,zhang2021rewardfree,sekar2020plan2explore,barreto2017successor}.
Typically exploration or representation learning ends before the downstream reward is given.
\taskname{} instead requires immediate execution from a frozen mid-action state, and its AND/OR/absence programs make utility nonlinear in primitive evidence.

\section{Baseline Adaptations}
\label{app:baselines}

Every baseline selects from exactly the same macro-actions.
Nearest Frontier uses path length to the closest reachable frontier.
Coverage and Information Gain score expected newly observed free space or entropy reduction.
Semantic Novelty counts new grounded categories under the shared detector.
ECC-style scoring estimates newly reached environment checkpoints; belief uncertainty prioritizes regions with uncertain spatial state; semantic-affordance priority ranks transit anchors; and latent-fact coverage counts newly supportable context facts.
Expected Requirement Deficit weighs missing primitives by the declared task prior.
Successor Features predicts discounted primitive occupancies and applies GPI for catalog weights.
Myopic Endpoint Utility uses the frozen executor to supervise terminal gain.
The reward-free model baseline plans for latent model disagreement.

Structural ablations retain training data and capacity whenever possible.
Endpoint-\jet{} removes all prefix labels; Prefix-Uniform retains prefixes but replaces $f_t$ with a uniform distribution; Independent-\jet{} uses a product of marginal hitting-time models; No-Certification admits raw detections; and Direct Utility predicts catalog-averaged utility without the evidence interface.
Task-Known receives $z$, Arrival-Known receives $\tau$, Oracle Memory receives true evidence state, Hindsight Candidate selects the best candidate after rollout, and Online Branch Planner may query simulator futures.
Only the first group constitutes fair task-hidden comparisons.

\section{Censored Evidence-Time Likelihood}
\label{app:likelihood}

For primitive $k$, let $s_k$ be its first certified prefix and $\delta_k=1$ if a hit occurs before $L_a$.
The right-censored negative log-likelihood is
\begin{equation}
\mathcal{L}_{\mathrm{hit}}=
-\sum_k\left[
\delta_k\log p_\theta(T_k^a=s_k)
+(1-\delta_k)\log p_\theta(T_k^a>L_a)
\right].
\end{equation}
For discrete hazards,
$p_\theta(T_k^a>L_a)=\prod_{r=1}^{L_a}(1-\alpha_{kr})$,
with mixture likelihood marginalized before taking the logarithm.
Calibration is evaluated by prefix horizon and primitive type; C-index measures hitting-time ordering under censoring.

\section{Atomic Interruption Protocol}
\label{app:atomic}

For every low-level step, the evaluator performs the following operations in order:
\begin{enumerate}[leftmargin=*]
  \item apply one control and advance the simulator;
  \item render the permitted agent observations;
  \item update pose belief and run the frozen perception stack;
  \item atomically commit newly validated records to \rfm{};
  \item draw or reveal the pre-sampled arrival event for the completed step;
  \item if the event occurs, cancel all remaining low-level controls and freeze memory, pose belief, and compute state.
\end{enumerate}
The post-arrival executor receives no observation generated after the freeze.
Arrival seeds are sampled independently of method behavior after conditioning on the specified hazard.
When a policy changes wall-clock inference time enough to matter, the benchmark offers both environment-step and wall-clock arrival tracks.

\section{Planning Algorithm}
\label{app:algorithm}

\begin{enumerate}[leftmargin=*,label=\textbf{\arabic*.}]
  \item Given memory $M_t$, pose belief $x_t$, task prior $p$, and arrival model
  $(f_t,S_t)$, construct
  $\mathcal{A}_t\gets\textsc{GenerateCandidates}(M_t,x_t)$.
  \item For every $a\in\mathcal{A}_t$, sample
  $\{\mathbf{T}^{a,(m)},\hat x^{a,(m)}_{1:L_a}\}_{m=1}^{P}
  \sim\jet_\theta(M_t,x_t,a)$.
  \item For each prefix $s$, compute
  \[
  \bar R(a,s)\gets \sum_z p(z)\frac{1}{P}\sum_m
  g_z(M_t\oplus\{k:T_k^{a,(m)}\le s\},\hat x_s^{a,(m)};B).
  \]
  \item Score each action by
  \[
  \hat Q(a)\gets
  \sum_s f_t(s)\bar R(a,s)+S_t(L_a)\hat V(M_{t+L_a},\hat x_{t+L_a})
  -\lambda C(a).
  \]
  \item Return $\arg\max_{a\in\mathcal{A}_t}\hat Q(a)$.
\end{enumerate}

\section{Evidence Schema and Requirement Programs}
\label{app:programs}

Each primitive has a type signature, a validator, a confidence calibration set, and a provenance requirement.
Grounded location requires a persistent entity identifier and 3D support.
An attribute view requires minimum projected size, visibility, and a calibrated attribute score.
A relation view requires compatible timestamps or geometrically registered observations.
Coverage certificates store observed free-space or surface support within a declared scope.
Absence combines coverage with category-specific detector recall; the system abstains when either component falls below threshold.
Reachability evidence records local-planner feasibility from the interrupted pose belief.

\begin{table*}[h]
  \centering
  \placeholdertable{3.2cm}{Full task-template inventory and program definitions: argument types, primitive requirements, utility, abstention rule, and composition split.}
  \caption{\textbf{Task program specification.}
  This table will be populated from the benchmark manifest so that the paper and released evaluator use identical definitions.}
  \label{tab:programs}
\end{table*}

\section{Implementation Details}
\label{app:implementation}

The final implementation report will include:
model width, depth, attention heads, latent mixture rank, maximum memory records, candidate count, prefix-bin width, particles per action, planning depth, leaf-value architecture, optimizer, learning-rate schedule, batch size, training updates, validator thresholds, all loss weights, hardware, simulator throughput, and wall-clock cost.
Hyperparameters are selected on validation scenes only.
The default configuration is frozen before any final test run.
All baselines use the same perception checkpoints and local controller.

\begin{table*}[h]
  \centering
  \placeholdertable{3.0cm}{Architecture, optimization, planning, memory, perception, navigation, hardware, and compute hyperparameters.}
  \caption{\textbf{Implementation configuration.}
  Every value is to be read from the released experiment configuration rather than transcribed manually.}
  \label{tab:hyperparameters}
\end{table*}

\section{Cross-Environment Protocol Details}
\label{app:cross-environment}

\paragraph{ProcTHOR.}
ProcTHOR supplies diverse interactive houses for counterfactual branch generation and held-out task compositions.
Its structured object, receptacle, attribute, and room metadata supports all four task families, but only information recovered through permitted egocentric observations may enter \rfm{} or the downstream executor.

\paragraph{HM3D and HM3D-OVON.}
HM3D is the primary photorealistic mesh domain and supports the complete navigation, attribute-QA, spatial-QA, and presence/absence catalog.
The cross-domain navigation slice uses the official HM3D-OVON unseen-scene split and open-vocabulary goal semantics.

\paragraph{OpenEQA.}
We use the active-exploration portion of OpenEQA and retain natural questions whose answers can be certified by the declared attribute, spatial-relation, or presence/absence evidence primitives.
The realized question is hidden before arrival; after interruption, evaluation follows the benchmark answer protocol with an additional abstention-aware report for questions lacking sufficient committed evidence.

\paragraph{HSSD ObjectNav.}
HSSD supplies high-quality, authored synthetic homes and semantically annotated object instances.
Its ObjectNav episodes test whether policies learned from larger procedural corpora transfer to more realistic scene composition while preserving the task-hidden arrival wrapper.

\paragraph{Habitat-GS.}
Habitat-GS PointNav tests whether prefix timing and planning remain effective when observations are rendered from 3D Gaussian splats.
It is a renderer-and-control stress test rather than the primary semantic benchmark.

\paragraph{Dynamic HSSD and P4D-Bench.}
The evolving-scene track uses versioned HSSD histories and P4D-Bench task construction.
Pre-task evidence is timestamped, the realized downstream query remains hidden until arrival, and the evaluator exposes no future scene version.
This track jointly stresses interrupted acquisition and stale evidence; we therefore report results both over all episodes and stratified by whether the queried state changed since its last certified observation.

\paragraph{Common controls.}
For every domain, task and arrival seeds are sampled independently after conditioning on the declared catalog and hazard.
Methods share the same scene version, start, observation stream, candidate generator, local controller, downstream executor, and service budget within a comparison.
Native SR, SPL, and QA metrics are retained alongside Recall-AWU and Service-AWU.
No aggregate score is computed across benchmark columns.

\section{Real-World Deployment}
\label{app:realworld}

\paragraph{Robot platforms.}
We deploy \method{} on the same physical configuration used by our companion robotic system: DEEP Robotics LYNX M20, X30, and Lite3.
The wheel-legged M20 is the primary quantitative platform; its standing size is $820\times430\times570$\,mm, mass is approximately 35\,kg, and manufacturer-rated endurance/range is 3\,h/15\,km unloaded.
It carries dual 96-line LiDARs and wide-angle cameras.
The quadrupedal X30 measures $1000\times695\times470$\,mm, has a mass of 56\,kg and a manufacturer-rated endurance/range of 2.5--4\,h/$\geq10$\,km, and provides integrated perception and autonomous navigation.
The LiDAR-equipped quadrupedal Lite3 measures $610\times370\times496$\,mm, has a mass of 13.5\,kg and a manufacturer-rated endurance/range of 1.5--2\,h/2.7\,km, and exposes obstacle avoidance and autonomous navigation through its platform interface.
Rated endurance and range are configuration descriptors, not experimental outcomes; time and distance are compared only between methods run on the same platform.

\paragraph{Deployment interface and scenes.}
Trials take place in multiple real indoor and outdoor environments.
Indoor trials include offices, corridors, meeting rooms, kitchens, and shared service areas with clutter, doors, furniture, objects, and people; outdoor trials include building approaches, paved walkways, courtyards, and open service routes with longer traversals and changing illumination.
The task catalog spans object or person navigation, visible-attribute questions, spatial-relation questions, and presence/absence verification.
Targets remain stationary within an individual post-arrival execution, while their placement and the surrounding context may differ across sessions.
Platform-specific perception and low-level navigation are abstracted behind a common macro-action interface, while \rfm{}, \jet{}, task programs, and the high-level arrival-aware planner remain unchanged.
M20 is used for the main matched comparison, while X30 and Lite3 test whether the high-level policy transfers across wheel-legged and quadrupedal embodiments.
On every platform, an arrival event cancels the remaining macro-action controls after the current observation is committed; downstream execution starts from that physical pose and frozen memory.
Native collision avoidance and emergency-stop facilities remain active throughout all trials.

\begin{table}[h]
  \centering
  \small
  \begin{tabular}{@{}llll@{}}
    \toprule
    \textbf{Platform} & \textbf{Morphology} & \textbf{Primary sensing} & \textbf{Evaluation role} \\
    \midrule
    LYNX M20 & Wheel-legged & Dual 96-line LiDAR + cameras & Main indoor/outdoor comparison \\
    X30 & Quadruped & Integrated perception stack & Cross-platform transfer \\
    Lite3 & Quadruped & LiDAR navigation stack & Lightweight transfer \\
    \bottomrule
  \end{tabular}
  \caption{\textbf{Physical platforms and roles.} Manufacturer specifications describe the tested configurations and are not treated as experimental outcomes.}
  \label{tab:robot-platforms}
\end{table}

\begin{table*}[h]
  \centering
  \placeholdertable{3.0cm}{Real-world results: method $\times$ indoor/outdoor Recall-AWU, Service-AWU, task success, distance, time, interrupted-prefix distribution, intervention rate, and 95\% confidence intervals; cross-platform transfer on M20, X30, and Lite3.}
  \caption{\textbf{Physical robot evaluation.}
  Populate every cell from completed trials under the matched real-world protocol.}
  \label{tab:realworld}
\end{table*}

\begin{figure*}[h]
  \centering
  \placeholderfigure{4.2cm}{Indoor and outdoor robot trajectories: pre-task exploration, task arrival inside a macro-action, committed evidence at interruption, and downstream task execution on M20, X30, and Lite3.}
  \caption{\textbf{Real-world deployments.}
  Qualitative frames will show the exact interruption prefix and evidence available when each request arrives.}
  \label{fig:realworld}
\end{figure*}

\section{Statistical Analysis}
\label{app:statistics}

The primary comparison is \method{} against the strongest validation-selected task-hidden baseline.
Training uses at least five independent seeds.
Confidence intervals use a hierarchical bootstrap that resamples scenes, episodes within scenes, and task/arrival draws within episodes.
Pairwise differences reuse identical scene, start, task, arrival, and noise seeds.
We report means, 95\% confidence intervals, and paired standardized effect sizes.
The test split is run once after thresholds, the robust radius, and early stopping are fixed.
Family utilities are normalized with validation-set constants and macro-averaged.
Per-family raw metrics remain visible to prevent a favorable aggregate from hiding a failed task family.

\section{Additional Result Placeholders}
\label{app:results}

\begin{table*}[h]
  \centering
  \placeholdertable{3.1cm}{Per-family navigation, attribute QA, spatial QA, and presence/absence results with success, SPL/cost, accuracy, grounding precision, risk--coverage, ECE, and Brier score.}
  \caption{\textbf{Per-task-family results.}}
  \label{tab:family}
\end{table*}

\begin{table*}[h]
  \centering
  \placeholdertable{3.1cm}{Ablations for endpoint prediction, prefix weighting, joint evidence, certification, utility consistency, ranking loss, and robust planning.}
  \caption{\textbf{Component ablations.}}
  \label{tab:component-ablations}
\end{table*}

\begin{figure}[H]
  \centering
  \placeholderfigure{4.5cm}{Qualitative trajectories and readiness curves: successful early evidence, endpoint trap, false certificate, candidate omission, and recovery after blockage.}
  \caption{\textbf{Qualitative cases and failure taxonomy.}}
  \label{fig:qualitative}
\end{figure}

\section{Reporting Checklist}
\label{app:checklist}

Before submission, the release and manuscript should verify:
\begin{itemize}[leftmargin=*]
  \item no realized task, target category, arrival sample, future frame, navmesh, or shortest-path distance leaks into pre-task action selection;
  \item interrupted actions do not execute to completion and the same atomic ordering is used by every method;
  \item all result placeholders are replaced by evaluator-generated artifacts and no draft values remain;
  \item all citations and venue metadata are manually verified;
  \item scene, task-composition, and arrival seeds are published with licenses;
  \item counterfactual branch count, simulator hours, training compute, inference latency, and peak memory are reported;
  \item physical logs record platform, environment, task and arrival seeds, interruption prefixes, interventions, and site/data permissions;
  \item confidence intervals reflect scene-level variation and all primary comparisons are paired;
  \item claims distinguish static-world, evolving-scene, and physical evaluations and remain restricted to the declared task families and benchmark protocols.
\end{itemize}
